\documentclass[10pt,a4paper]{article}
% 815agent 工程问答 16 题 — 由 interview.html 整理, XeLaTeX 中文排版
% 小五 (9pt) 正文字号
\usepackage{fontspec}
\usepackage[UTF8,fontset=none]{ctex}
\setCJKmainfont{Microsoft YaHei}[BoldFont={Microsoft YaHei Bold}]
\setCJKsansfont{Microsoft YaHei}[BoldFont={Microsoft YaHei Bold}]
\setCJKmonofont{Microsoft YaHei}[BoldFont={Microsoft YaHei Bold}]
\setmainfont{Arial}
\setsansfont{Arial}
\setmonofont{Consolas}

\usepackage[a4paper,margin=1.6cm,top=1.4cm,bottom=1.4cm]{geometry}
\usepackage{xcolor}
\usepackage{colortbl}
\usepackage{longtable}
\usepackage{array}
\usepackage{enumitem}
\usepackage{titlesec}
\usepackage{fancyhdr}
\usepackage[hidelinks]{hyperref}

% ---------- 配色(对应 HTML 深色页的语义色,转浅底) ----------
\definecolor{acc}{HTML}{0B62A2}     % 主蓝
\definecolor{okc}{HTML}{1B7F4B}     % 已有机制·绿
\definecolor{addc}{HTML}{0B62A2}    % 针对性补全·蓝
\definecolor{dimc}{HTML}{5A6672}    % 参照·灰
\definecolor{linec}{HTML}{B8C2CC}   % 表格线
\definecolor{headc}{HTML}{EAF1F8}   % 表头底
\definecolor{boxbg}{HTML}{F5F8FB}   % 问答卡底

% ---------- 小五字号:9pt ----------
\renewcommand{\normalsize}{\fontsize{9}{13.2}\selectfont}
\renewcommand{\small}{\fontsize{8}{11.5}\selectfont}
\renewcommand{\footnotesize}{\fontsize{7.5}{10.5}\selectfont}
\renewcommand{\large}{\fontsize{12}{16}\selectfont}
\renewcommand{\Large}{\fontsize{15}{20}\selectfont}

% ---------- 标题 ----------
\titleformat{\section}{\Large\bfseries\sffamily\color{acc}}{\thesection}{0.6em}{}
[\vspace{-0.4em}{\color{acc}\rule{\textwidth}{1.2pt}}]
\titlespacing*{\section}{0pt}{12pt}{8pt}
\setcounter{secnumdepth}{0}

% ---------- 页眉页脚 ----------
\pagestyle{fancy}
\fancyhf{}
\fancyhead[L]{\footnotesize\sffamily\color{dimc}815agent 实战报告 · 工程问答(面试题提取 × 16 题)}
\fancyhead[R]{\footnotesize\sffamily\color{dimc}\thepage}
\renewcommand{\headrulewidth}{0.4pt}
\renewcommand{\headrule}{\color{linec}\hrule width\headwidth height\headrulewidth}

% ---------- 标签 ----------
\newcommand{\taghave}{\textcolor{okc}{\textbf{〔已有机制〕}}}
\newcommand{\tagadded}{\textcolor{addc}{\textbf{〔针对性补全〕}}}
\newcommand{\tagref}[1]{\textcolor{dimc}{\footnotesize\{#1\}}}

% ---------- 问答卡 ----------
% 题号小色块 + 加粗题目一行, 正文平铺(不用嵌套盒, 兼容 Q11 longtable)
\newcommand{\qhead}[2]{\par\vspace{7pt}\noindent
  \colorbox{acc}{\color{white}\bfseries\small #1}\hspace{4pt}\textbf{#2}\par\vspace{2pt}}
\newenvironment{qa}[2]{\qhead{#1}{#2}\small}{\par\vspace{2pt}}

\setlength{\parskip}{2pt}
\setlength{\parindent}{0pt}

\begin{document}

\begin{center}
  {\Large\bfseries\sffamily 工程问答 --- 面试题提取 $\times$ 815agent 实战回答}\\[4pt]
  {\small\color{dimc}题目来自 DeepSeek Agent 岗三面连环追问(工程细节/系统设计/赛道认知)。共 16 题,每题标注:
  \textcolor{okc}{\textbf{〔已有机制〕}}现场验证过 \quad
  \textcolor{addc}{\textbf{〔针对性补全〕}}本轮新增 \quad
  \textcolor{dimc}{\{kimi/Hermes 参照\}}}
\end{center}
\vspace{-4pt}

% ================= 一面 =================
\section{一面 · 工程细节题}

\begin{qa}{Q1}{上下文用的是 message window 还是 token window?线上为什么这么选?}
\taghave{} \textbf{token window}。\texttt{context.py} 以估算 token 数(chars/4)对 30 万窗口设 85\% 触发线,message window 无法表达``10 条消息里有一条 240KB 的工具输出''这种真实分布---T5 实测单条 read\_file 结果就是 25 万字符。选 token window 的代价是估算漂移,工业解法是 kimi 的``实测锚点+估算尾部''双层计量,815agent 停在单层粗估 + 85\% 保守阈值 + 溢出自愈兜底(Q4)。\tagref{kimi: 双层计量 · Hermes: 三套信号协调}
\end{qa}

\begin{qa}{Q2}{会话持久化的过期策略怎么定?遇到过并发会话冲突吗?}
\tagadded{} 存储选型是 \textbf{JSONL 文件而非 Redis}(教程准则四:只要恢复$\to$追加日志够用,要上检索才换库)。过期策略:\texttt{session.cleanup\_old\_sessions(workdir, ttl\_days=30)},按文件 mtime TTL 清理,CLI \texttt{-{}-cleanup-days N}。\\
并发冲突遇到过,分两处解决:① 多会话同 workdir 写 MEMORY.md $\to$ 行级原子追加 + tmp+rename 原子替换截断;② 多子代理共享上下文压缩 $\to$ 物理隔离(各自独立 transcript),参考 Hermes \#38727 的尸检。\tagref{kimi: AppendLogStore tmp+rename cutover · Hermes: SQLite 压缩锁 TTL 300s}
\end{qa}

\begin{qa}{Q3}{长期记忆用向量库吗?记忆越来越多、检索越来越不准怎么办?}
\tagadded{} \textbf{不用向量库}---815agent 的记忆是有 8000 字符硬上限的平文件(容量即精度:上限强制淘汰旧条目,检索域永远小到关键词足够准)。本轮补了 \texttt{memory\_search} 工具(关键词检索,只读)。真正的防线在写入侧:provenance 时间戳溯源 + 注入安检,毒条目不进检索域。什么时候该上向量库:记忆条目破千、跨项目复用时---教程里 Hermes 用 SQLite FTS5(含 CJK trigram)也是先 BM25 后向量,且曾砍掉摘要 LLM 路径回归纯 BM25。\tagref{Hermes: FTS5+trigram,``session 霸榜''教训}
\end{qa}

\begin{qa}{Q4}{上下文 token 溢出的兜底方案是什么?}
\tagadded{} 三层:① 预防---每步 preflight 估算超 85\% 即压缩;② \textbf{自愈(本轮新增)}---provider 报 context/overflow/length 错误时 \texttt{\_chat()} 捕获,\texttt{force=True} 强制压缩后重试,连续 3 次放弃(对齐 kimi 溢出自愈上限);③ 兜死---摘要调用失败 fail-open 保原会话,预算耗尽还有 grace 收尾。单测 \texttt{TestOverflowHeal} $\times$2 覆盖成功与放弃两条路径。\tagref{kimi: 错误处理器认领 413,下调有效窗口 $\times$0.85}
\end{qa}

\begin{qa}{Q5}{单用户日均 token 成本多少?百万级流量怎么控制?}
\tagadded{} 可测才能答:本轮给每步 LLM 调用加了\textbf{用量账本}(prompt/completion/calls 累加,turn 结束落 transcript usage 事件)。现场实测:一次``建文件+运行+查记忆''的轻任务 = \textbf{5,838 prompt + 189 completion / 4 次调用};一次完整项目任务(T1,39 测试项目)约 14 次工具调用、$\sim$10 次 LLM 调用。百万级控制三板斧全在架构里:① 缓存宪法(前缀只追加,DashScope 命中缓存约 1/10 价);② 结果治理(50K 落盘只回 2K,T5 实测 25.6 万字符$\to$2,130);③ 预算硬墙(max\_steps + 子代理独立预算)。\tagref{kimi: 分段缓存+末尾注入 · Hermes: 4 断点+冻结快照}
\end{qa}

\begin{qa}{Q6}{调用拦截:用户恶意刷长文本怎么限制?}
\tagadded{} 三层拦截:① 入口---\texttt{AGENT815\_MAX\_INPUT\_CHARS}(默认 20 万字符)超限直接拒绝,\textbf{不消耗任何 LLM 调用}(单测验证 FakeLLM 零调用);② 工具侧---bash 120s 超时、输出 30K 截断、大结果 50K 落盘;③ 循环侧---max\_steps 硬墙 + stop hook 一次性闩锁防``永不结束''。
\end{qa}

\begin{qa}{Q7}{为什么用基座模型不用微调?拿什么数据做判断?}
\taghave{} 选基座(qwen3.7-max)因为 Harness 架构把行为约束都放在了\textbf{工程层}(system prompt 纪律 + 工具协议 + 权限链),这些不需要权重级定制;微调的适用面是输出格式/领域风格稳定且高频的场景,coding agent 的需求是通用推理 + 工具调用可靠性。判断数据:① 工具调用成功率(实测 73 次调用 0 协议错误);② 任务完成率(8 个现场任务 8/8 exit 0);③ 失败自愈率(T1/T3 测试失败$\to$edit\_file 修复$\to$复测通过);④ 单任务 token 成本(Q5 账本)。微调唯一可能值回票价的点是把 BASE\_PROMPT 的纪律烧进权重省 system prompt token---但会损失迭代速度,不划算。
\end{qa}

\begin{qa}{Q8}{讲一个线上故障排查:改了哪块逻辑,效果多少?}
\taghave{} \textbf{单 user 会话压缩失效}。现象:小窗口压力测试(MAX\_CONTEXT=1500)下压缩事件恒为 0。排查:逐步重放真实 transcript 计算各 preflight 时点估算,确认第 3 条消息起已超阈值却不触发 $\to$ 定位到 \texttt{context.py} ``最后一条 user 必须在 tail''守护规则:一次性任务唯一 user 消息在 head(索引 1),守护把 cut 强制拉回 1,middle 恒空,压缩静默放弃。修复:守护仅在 \texttt{last\_user\_idx $\geq$ 2} 生效。效果:同任务压缩触发从 \textbf{0 $\to$ 2 次},免疫摘要正确注入;回归测试 \texttt{test\_single\_user\_message\_still\_compacts} 防复发。教训:守护规则必须考虑``被保护对象已在保护区''的边界。
\end{qa}

% ================= 二面 =================
\section{二面 · 系统设计题}

\begin{qa}{Q9}{系统模块怎么拆分?职责与通信?为什么这么拆?}
\taghave{} 11 个模块按``窄腰核心 + 边上插拔''拆(见模块设计页):loop 是唯一有状态核心,其余模块通过三个窄接口通信---① 工具协议(schema+handler 字典);② 事件表(hooks 5 事件);③ 消息列表(唯一共享状态,只追加)。拆分依据是教程两条宪法:核心小到可以穷举测试(loop 168 行),换 provider/换存储/换表面不动核心。通信不走消息总线而共享不可变前缀的消息列表,是为了缓存宪法---任何总线式动态注入都会打碎前缀缓存。
\end{qa}

\begin{qa}{Q10}{怎么判断单工具简单任务 vs 多步拆解复杂任务?拆解粒度?}
\taghave{} 815agent 走 kimi 路线:\textbf{拆解权交给模型,工程层只设边界}---模型自主决定是否连续调工具(终止由模型决定),工程层的控制是:max\_steps 硬墙(60)、子代理独立预算(30)、\texttt{delegate\_task} 让模型把``探索性子任务''显式派生出去(粒度信号:会污染主线上下文的活就派生,探索日志只回蒸馏报告)。Hermes 路线是预算 refund(模型写代码循环调工具不侵蚀预算)---两种都验证了``不要在工程层替模型做拆解判断'',那会变成永远追不上模型能力的启发式。
\end{qa}

\begin{qa}{Q11}{任务执行一半失败:重试、回滚还是重新规划?触发条件?}
分四层,每层触发条件显式:

{\small
\begin{longtable}{|>{\raggedright\arraybackslash}p{0.16\textwidth}|>{\raggedright\arraybackslash}p{0.46\textwidth}|>{\raggedright\arraybackslash}p{0.30\textwidth}|}
\hline
\rowcolor{headc}\textbf{失败类型} & \textbf{策略} & \textbf{触发条件}\\\hline
\endhead
工具执行失败 & 错误文本回填,\textbf{模型自行重试或换路}(T1 实测:测试失败$\to$edit\_file$\to$复测通过) & 任何 tool exception\\\hline
429/5xx/网络超时 & 指数退避重试($\times$2 上限 30s + 20\% 抖动 + Retry-After),最多 5 次 & HTTP 状态码/连接异常\\\hline
上下文溢出 & 强制压缩后重试,连续 3 次放弃 & 错误文本含 context/overflow/length\\\hline
权限拒绝/hook 阻断 & \textbf{不重试不绕过},原因回填让模型换方案(T4b 实测 pip 被禁$\to$标准库完成) & sandbox/hook veto\\\hline
\end{longtable}}

回滚在最小设计里刻意没有(文件操作无事务),对应演化级是 kimi 的 undo/wire Op 重放---痛了再加。
\end{qa}

\begin{qa}{Q12}{多工具怎么联动?权限怎么管控?怎么防越权?}
\taghave{} 联动:工具是纯字典注册表,联动顺序由模型按 schema 描述编排;结果治理保证任何单工具输出不灌爆窗口。权限五道防线,顺序即优先级:① \textbf{hardline 硬线}(rm -rf /、fork bomb、写 authorized\_keys,\textbf{yolo 也不放行}---用户不能授权无恢复路径的破坏);② 去混淆(NFKC 全角/\$IFS/引号拼接)防绕过;③ 未知工具 fail-closed;④ 只读/写执行三分级 + ask/auto/yolo 模式;⑤ workdir 围栏 + 子进程环境变量剥离(KEY/TOKEN 类,实测工作区零密钥痕迹)。防越权的价值观:\textbf{危险检测先于用户授权}。\tagref{Hermes: hardline 先于 yolo bypass · GHSA-rhgp 凭据隔离}
\end{qa}

% ================= 三面 =================
\section{三面 · 认知题(回答框架,以本项目为论据)}

\begin{qa}{Q13}{Agent 是长期风口还是短期热点?长期落地方向?}
长期。论据:模型能力每代跃迁,但教程里两个工业 codebase 的厚度证明\textbf{工程表现的上下限由 Harness 决定}---同一模型接不同 Harness,安全性/成本/长任务存活率差一个数量级,这层工程不会随模型变强而消失(缓存宪法、权限纵深、状态可恢复都是模型外问题)。落地方向:有人值守的编程 CLI(kimi 型)先跑通,无人值守的多表面平台(Hermes 型)跟随,自学习闭环(记忆/技能双闭环)是拉开差距的下半场。
\end{qa}

\begin{qa}{Q14}{和普通大模型开发者比,核心竞争力?}
能从``事故''反推架构。815agent 每个模块都能回答``防什么事故、对应演化阶梯第几级、kimi/Hermes 怎么解''---demo 开发者堆框架,工程师管爆炸半径(hardline 先于 yolo)、管成本(缓存宪法)、管最坏情况下限(fail-open 方向都是显式决策)。
\end{qa}

\begin{qa}{Q15}{1-2 年内最大瓶颈是模型能力还是工程能力?ToB 还是 ToC 先跑通?}
工程能力。模型已经``够用''(实测 qwen3.7-max 零协议错误完成 8 个完整项目),卡脖子的是:长会话成本(压缩与缓存的博弈)、无人值守安全(纵深防御)、状态可恢复(断点续跑)。ToB 先跑通---容错预算高、场景边界清晰、权限模型可预先声明;ToC 的越权风险面不可控。
\end{qa}

\begin{qa}{Q16}{从 0-1 带一条 Agent 业务线,第一步做什么?}
先定形态再写代码(教程决策树 Q1):有人盯着还是无人值守?这一个问题决定主循环形状(插件链 vs 预算硬墙)、安全厚度(规则链 vs 纵深栈)、存储选型(JSONL vs SQLite+FTS5)。招人按模块招(循环/上下文/安全/记忆技能/表面),每人负责一条演化阶梯的爬升路线,评审标准就是``每份复杂性答得出防什么事故''。
\end{qa}

\vspace{6pt}
{\footnotesize\color{dimc}\rule{\textwidth}{0.4pt}\\
本页为 815agent 实战报告独立附页(源:interview.html)。完整版见 815agent\_report\_full.pdf(六大板块合并)。}

\end{document}
